\documentclass{beamer}
\usepackage{graphicx}
\usepackage{amssymb}
\usepackage{amsmath}
\usepackage{amsthm}
\usepackage{bm}
\usepackage{color}
\usepackage{array}
\usepackage{enumitem}

%%%%%%%%%%%%%%%%
\newcommand{\rhoo}{\rho\mkern-2mu_{\mathsf{\scriptscriptstyle o}}}
\newcommand{\ho}{h_{\mathsf{\scriptscriptstyle o}}}
\newcommand{\vo}{v_{\mathsf{\scriptscriptstyle o}}}
\newcommand{\nablah}{\nabla\mkern-5mu_{\mathsf{\scriptscriptstyle h}}}
%%%%%%%%%%%%%%%%


%%citations%%%%%%%%%%%%%%%%%%%%%%%%%%%%%%%%%%%%%%%%%%%%%%%%%%%%%% 
\setbeamerfont{bibliography entry note}{size=\tiny}
\usepackage[T1]{fontenc}    % Fixes the underscore/dot and French accents

\usepackage[
    backend=biber,
    style=ext-authoryear,
    articlein=false,    % <--- This specifically removes "In:" for @article
    % --- IN-TEXT CITATION SETTINGS ---
    maxcitenames=2,   % Use "et al." if there are more than 2 authors
    mincitenames=1,   % When truncating, show only the 1st author
    % --- BIBLIOGRAPHY (SLIDE) SETTINGS ---
    maxbibnames=99,   % Show up to 99 authors (essentially all)
    minbibnames=1,    % Fallback if a list is somehow longer than 99
    backref=true,      % <--- Enables "back to text" links
    uniquelist=false,  % Keeps "et al." consistent even with similar entries
    giveninits=true,   % <--- This produces "Griffies, S. M."
]{biblatex}

\renewcommand*{\bibfont}{\footnotesize} % Or \scriptsize if you're desperate
\DefineBibliographyStrings{english}{%
  backrefpage = {cited on page},
  backrefpages = {cited on pages},
}

\setlength{\bibitemsep}{0pt}  % Space between different entries
\setlength{\bibparsep}{0pt}   % Space between paragraphs within an entry


%\bibliography{Griffies_book_refs}
\addbibresource{./../../../../../Princeton_University/courses/Princeton_AOS/GFM_book/Griffies_book_refs.bib}


% --- BEAMER SPECIFIC CLEANUP ---
% This removes the "document/article" icon in the bibliography
\setbeamertemplate{bibliography item}{}

% This makes the font size smaller for bibliography slides
\setbeamerfont{bibliography entry author}{size=\scriptsize}
\setbeamerfont{bibliography entry title}{size=\scriptsize}
\setbeamerfont{bibliography entry location}{size=\scriptsize}
\setbeamerfont{bibliography entry note}{size=\scriptsize}

%%%%%%%%%%%%%%%%%%%%%%%%%%%%%%%%%%%%%%%%%%%%%%%%%%%%%%%% 

\hypersetup{
    colorlinks=true,   % Colors the text instead of drawing a box
    citecolor=blue,    % Makes "Haine et al 2025" blue and clickable
    linkcolor=black,   % Keeps internal links (like TOC) black and clean
    urlcolor=blue      % Keeps web links blue
}

\definecolor{darkblue}{rgb}{0,0,0.5} % Professional navy for your slides


%
% Choose how your presentation looks.
%
% For more themes, color themes and font themes, see:
% http://deic.uab.es/~iblanes/beamer_gallery/index_by_theme.html
%
\mode<presentation>
{
  \usetheme{default}      % or try Darmstadt, Madrid, Warsaw, ...
  \usecolortheme{default} % or try albatross, beaver, crane, ...
  \usefonttheme{default}  % or try serif, structurebold, ...
  \setbeamertemplate{navigation symbols}{}
  \setbeamertemplate{caption}[numbered]
  \setbeamertemplate{footline}[frame number]
} 

\usepackage[english]{babel}
\usepackage[utf8x]{inputenc}


% to center the title to each slide 
\setbeamertemplate{frametitle}{
    \begin{centering}
        \insertframetitle\par
        \insertframesubtitle\par
    \end{centering}
}


% Sets the "Figure XX" part
\setbeamerfont{caption name}{size=\scriptsize}
% Sets the actual text of the caption
\setbeamerfont{caption}{size=\scriptsize}

\usepackage[skip=1pt]{caption}

\newcommand{\term}[1]{{\rmfamily\scshape #1}}

\setbeamerfont{frametitle}{size=\normalsize, series=\mdseries, family=\sffamily}

\AtBeginSection[]
{
  \begin{frame}{Outline}
    \tableofcontents[currentsection]
  \end{frame}
}

%%%%%%%%%%%%%%%%%%%%%%%%%%%%%%%%%%%%%%%%%%%%%%%%%%%%%%%%%%%%%%%%%%%%%%%%%%%

\title{A primer on the ocean mesoscale and its parameterization in
  circulation models \\
  \vspace{.2cm}
  \small Sampling the physics/maths/numerics (with a focus on tracers)}
\author{\small stephen.griffies@locean.ipsl.fr \\ 
  smg@princeton.edu \\
  stephen.m.griffies@gmail.com}
\institute{CNRS IPSL/LOCEAN + Princeton University}
\date{Talk given 18 May 2026 at the Institut Henri {Poincar\'{e}} workshop: \\ {\it Instabilities and Transitions in Geophysical Flows}

\vspace{1cm}
\centering
\includegraphics[width=.5\textwidth]{FIGS/ParisPrincetonLogos.pdf}
  
}


\begin{document}

\small 

\begin{frame}
  \titlepage
\end{frame}

% --- Table of Contents Slide ---
\begin{frame}{Outline}
  \tableofcontents
\end{frame}

\section{Introduction and motivation}

\begin{frame}{Audiences for this talk}

\begin{itemize}

\item[$\star$] Mathematical analysts interested in the physics of
  subgrid scale parameterizations used in ocean circulation models:
  \begin{itemize}
  \item[$\bullet$] {\bf Physics informing the maths.}
  \item[$\bullet$] For example, see \textcite{KornTiti2024}. 
  \end{itemize}
  
\item[$\star$] Physical analysts interested in the mathematical
  expression of subgrid scale operators used in circulation models:
    \begin{itemize}
     \item[$\bullet$] {\bf Maths informing the physics.}
    \end{itemize}

  \item[$\star$] Numerical analysts interested in realizing subgrid
    scale operators in circulation models.
    \begin{itemize}
     \item[$\bullet$] {\bf Numerics realizing the maths and physics.}
    \end{itemize}

\end{itemize}


\end{frame}


\begin{frame}{Ocean mesoscale focus of this talk}

  
\begin{figure}
\centering
\includegraphics[width=.6\textwidth]{FIGS/SUOMI_GulfStream.pdf}
%\includegraphics[width=.4\textwidth]{FIGS/Southern_ocean_speed.jpg}
%\includegraphics[width=.3\textwidth]{FIGS/CM2.6_SPEED_ICE_2000528_NA_MILLER.png}

%\caption{\scriptsize Ocean surface speed from a global climate model
%  (GFDL/CM2.6) with roughly 5~km horizontal grid spacing in the
%  Southern Ocean \parencite{Griffies_etal_CM2p6_heat}.}
\caption{\scriptsize Surface image in North Atlantic from the
  NASA-NOAA SUOMI-NPP weather satellite, with Gulf Stream and
  associated mesoscale eddies, along with submesoscale fronts and
  eddies, highlighted by phytoplankton pigments. Image available from
  https://www.nasa.gov/image-article/turbulent-north-atlantic/.}
\end{figure}

We here focus on the ocean mesoscale, which is characterized by
     \begin{itemize}
       \footnotesize
     \item[$\bullet$] Small Rossby number (rotationally dominant)
     \item[$\bullet$] Large Richardson number (strong vertical stratification)  
     \item[$\bullet$] Quasi-geostrophy (flow scales $L \approx
       \mathcal{O}(L_{d})$) and planetary geostrophy (flow scales $L
       \gg \mathcal{O}(L_{d})$) inform the parameterizations ($L_{d} =
       N \, H/f$ is the deformation scale).
    \end{itemize}

\end{frame}



\begin{frame}{Perspective (or perhaps, an apology)}

\begin{itemize}

\item[$\star$] Mesoscale eddy parameterizations (and by extension,
  ocean eddies) are central to the ocean's role in the climate system:
  \begin{itemize}
    \footnotesize
  \item[$\bullet$] Ocean heat uptake is affected by turbulent motions,
    with $\approx 90\%$ anthropogenic heat absorbed by the ocean.
  \item[$\bullet$] Ocean carbon uptake is also affected by turbulent
    motion, with ocean accounting for $\approx 25\%$ of carbon budget.
  \item[$\bullet$] Sea level rise (thermosteric expansion + melting ice shelves) 
  \item[$\bullet$] Climate mean, variability, stability, and
    sensitivity to $CO_2$ increases.
  \end{itemize}

  
\item[$\star$] Practical problems related to ocean eddy
  parameterizations occupy the attention of many ocean theorists and
  model developers.
    \begin{itemize}
  \footnotesize
   \item[$\bullet$] Eddy parameterizations aim to encapsulate physics
     in a closed mathematical form, allowing for mathematical analysis
     and numerical experimentation.
   \item[$\bullet$] There is a trend in the community towards
     efficient and accurate dynamical cores, coupled to better
     computer hardware, that aspire to resolve ocean mesoscale eddies,
     rather than parameterize.
    \item[$\bullet$] Even so, parameterizations offer an effective way
      to engage a broad suite of scientists and mathematicians into
      the study of ocean circulation and the tole of ocean turbulence
      in climate.
  \end{itemize}

\end{itemize}


\end{frame}




\begin{frame}{How this talk fits into the workshop theme}

\begin{itemize}

\item[$\star$] Ocean turbulence is initiated by instabilities, such as
  baroclinic, symmetric, and Kelvin-Helmholtz (shear).

\item[$\star$] This talk is mostly concerned with parameterizing the
  effects from {\bf ocean mesoscale turbulence} (geostrophic
  turbulence) in a rotating and stratified ocean.

\item[$\star$] Ocean mesoscale turbulence is directly related to
  baroclinic instability, with secondary instabilities arising from
  shear instability.

\end{itemize}

\begin{figure}
\centering
\includegraphics[width=\textwidth]{FIGS/instabilityGiants.pdf}
\caption{\scriptsize Pioneers of fluid instabilities central to the varieties of
  ocean turbulence.}
\end{figure}


\end{frame}


\section{Physical phenomenology of ocean mesoscale flows}


\begin{frame}{Ocean circulation: horizontal and vertical}


\begin{figure}
\centering
\includegraphics[width=\textwidth]{FIGS/circulation_images.pdf}
\caption{\scriptsize Left: schematic of the horizontal ocean
  circulation patterns from Lynn Talley. Center: meridional-vertical
  overturning circulation, with Antarctica at the center
  \parencite{Talley2013}.  Right: neutral density (buoyancy) section
  in Pacific \parencite{Wunsch_Ferrari2004}.}
\end{figure}

Selected features of the ocean circulation:
      \begin{itemize}
      \scriptsize
      \item[$\star$] Broad horizontal gyre circulations coupled to the
      Antarctic Circumpolar Circulation.
    \item[$\star$] Overturning cells, connecting in the deep Southern
      Ocean and with the North Atlantic having the only deep
      penetration in the northern hemisphere.
    \item[$\star$] Buoyancy is characterized by light water above
      heavy, which provides a resistance to vertical motion and
      affects properties of ventilation (vertical transfer of
      tracers), waves, and eddies.
      \end{itemize}


\end{frame}



\begin{frame}{Multiscale pathways of ocean flows}

\vspace{-.3cm}
\begin{figure}
\centering
\includegraphics[width=\textwidth]{FIGS/Pathways-Roy.pdf}
\caption{\scriptsize Multiscale pathways of ocean flows from the
  review by \textcite{Aluie_etal2026}.}
\end{figure}

\begin{itemize}

  \item[$\star$] The mesoscales are among a wide suite of multi-scale
    ocean motions, with a transfer of mechanical energy between these
    motions indicative of coupling across scales.

\item[$\star$] Developing a theoretical and phenomenological
  understanding of these flows, and their interactions, is a goal of
  ocean physics.
  
\end{itemize}

\end{frame}



\begin{frame}{Ocean energetics across scales}


\begin{figure}
\centering
\includegraphics[width=\textwidth]{FIGS/ARFM_spectra.pdf}
\caption{\scriptsize Kinetic energy spectra of surface ocean
  geostrophic flow in the global ocean, computed from AVISO
  (satellite) and GLORYS (reanalysis) via coarse-graining:
  \textcite{Storer_etal2022,Storer_etal2023,Buzzicotti_etal2023,Aluie_etal2026}.}
\end{figure}


\vspace{-.3cm}
\begin{figure}
\centering
\includegraphics[width=\textwidth]{FIGS/Storer_etal_cascade_map.pdf}
\caption{\scriptsize Kinetic energy transfer from the GLORYS ocean
  product ($1/12^{\circ}$ NEMO model), from
  \textcite{Storer_etal2022}.  Blue means energy transfer to large
  scales, and red to small scales.}
\end{figure} 

\end{frame}





\section{Mathematical description of eddy tracer transport}

\begin{frame}{Dynamical core + subgrid scale parameterizations}

 Equations of an ocean model:
\begin{subequations}
 \begin{align}
   \frac{\mathrm{D}{\bm v}}{\mathrm{D}t}
   +
   2 \, {\Omega} \times {\bm v}
   &=
   -\nabla \Phi
   +
   \rho^{-1} \, (-\nabla p +  \color{red} \nabla \cdot \mathbb{T} \color{black})
  \\
  \color{red} \nabla \cdot \mathbb{T} \color{black}
  &= \text{\footnotesize divergence of subgrid stresses} 
    \\
  \rho \, \frac{\mathrm{D}C}{\mathrm{D}t} &= \color{red}  -\nabla \cdot {\bm J} \color{black} 
  \\
  \color{red} -\nabla \cdot {\bm J} \color{black}
  &= \text{\footnotesize convergence of subgrid tracer fluxes} 
 \end{align}
\end{subequations}

\begin{itemize}
\item[$\star$] Dynamical core
  \begin{itemize}
    \footnotesize
  \item[$\bullet$] That part of the model concerned with numerically
    representing the flow on the discrete grid (material time changes,
    Coriolis, gravity, pressure gradient).
  \end{itemize}
\item[$\star$] \color{red} Subgrid scale (eddy) parameterizations
  \color{black}
  \begin{itemize}
     \footnotesize 
   \item[$\bullet$] Closure theories that aim to parameterize impacts
     on the resolved flow from processes that live at scales smaller
     than the discrete grid.
     \item[$\bullet$] In a circulation model, the subgrid scale is
       operationally defined by the grid size.
  \end{itemize}
\end{itemize}

\end{frame}


\begin{frame}{Advection-diffusion equation for source-free (conservative) tracers}

\begin{itemize}

\item[$\star$] Boussinesq\footnote{The compressible (non-Boussinesq)
equations require only modest conceptual generalizations beyond the
Boussinesq ocean considered here.  Note that $\rhoo$ is the prescribed
reference density.} tracer advection-diffusion equation
\begin{equation}
    \partial_{t} C  + \nabla \cdot ( {\bm v} \, C)
    = -\nabla \cdot {\bm F} \quad \mbox{and} \quad \nabla \cdot {\bm v} = 0
    \quad \mbox{and} \quad {\bm J} = \rhoo \, {\bm F}.
\end{equation}

\item[$\star$] Subgrid scale (SGS) flux vector, ${\bm F}$, is written
  in terms of a second order transport tensor, $\mathbb{J}$, projected
  onto the tracer gradient
\begin{equation}
  {\bm F}  = -\mathbb{J} \cdot \nabla C
  \Longrightarrow
  -\nabla \cdot {\bm F} = \nabla \cdot (\mathbb{J} \cdot \nabla C). 
\end{equation}

\item[$\star$] Decompose the subgrid tracer transport tensor into two
  components: (i) positive and symmetric tensor, $\mathbb{K}$, and
  (ii) anti(skew)-symmetric tensor, $\mathbb{A}$
  \begin{subequations}
    \begin{alignat}{2}
    \mathbb{J} &=  \tfrac{1}{2} (\mathbb{J} + \mathbb{J}^{T})
     + \tfrac{1}{2} (\mathbb{J} - \mathbb{J}^{T})
     \quad &&\mbox{decompose} 
     \\
     \mathbb{K} &\equiv \tfrac{1}{2} (\mathbb{J} + \mathbb{J}^{T}) =  \mathbb{K}^{T}
     \quad &&\mbox{diffusion}
     \\ 
     \mathbb{A} &\equiv \tfrac{1}{2} (\mathbb{J} - \mathbb{J}^{T}) = -\mathbb{A}^{T}
          \quad &&\mbox{skew-diffusion / advection.}
    \end{alignat}
  \end{subequations}

\end{itemize}

\end{frame}



\begin{frame}{Distinguishing passive tracers and active tracers}

\begin{itemize}

\item[$\star$] {\bf Dynamically passive tracers}:
  concentration, $C$.
  \begin{itemize}
   \item[$\bullet$] Advection-diffusion equation is linear, with fluid
     flow (${\bm v}$) and mixing tensor ($\mathbb{J}$) independent of
     the tracer concentration, $C$.
    \item[$\bullet$] Admits a Green's function solution
      \parencite{Haine_etal2025}, for use in diagnosing tracer
      timescales, ventilation pathways, ideal age.
   \end{itemize}
  

\item[$\star$] {\bf Dynamically active tracers}: (Conservative
  Temperature, $\Theta$, and salinity, $S$)
  \begin{itemize}
  \item[$\bullet$] Advection-diffusion equation is nonlinear, with the
    flow and the mixing tensor a function of temperature and salinity.
  \item[$\bullet$] Special care is needed when numerically realizing
    the subgrid scale transport tensor, to avoid nonlinear
    instabilities such as those identified by
    \textcite{GriffiesIsoneutral}.
  \item [$\bullet$] The associated mathematical well-posedness
    analysis is particularly tricky \parencite{KornTiti2024}.
  \end{itemize}
  
\end{itemize}

\end{frame}


\begin{frame}{Two basic types of eddy tracer transport: stirring and mixing}

The symmetric/anti-symmetric decomposition of the tracer transport
tensor reflects two kinematically distinct forms of eddy tracer
transport.

\begin{figure}  
\centering
\includegraphics[width=.65\textwidth]{FIGS/Welander_stir_mix.pdf}
\caption{\scriptsize Schematic from \textcite{Welander1955}
  illustrating the enhanced gradients from advection, thus exposing
  tracer concentration to mixing from molecular diffusion. See also
  \textcite{Eckart1948,MullerGarrett2002,Garrett2006}.}
\end{figure}


\begin{itemize}
  %\scriptsize 
  \item[$\bullet$] $\mathbb{K}$: diffusion tensor, parameterizing
    subgrid processes leading to irreversible downscale mixing that
    reduces tracer gradients (dissipation).
  \item[$\bullet$] $\mathbb{A}$: skew-diffusion tensor, parameterizing
    subgrid processes leading to advective fluxes that stir tracers
    that can increase gradients.
  \item [$\bullet$] Eddy-advection is related to Stokes' drift
    \parencite{Moffatt1983,MiddletonLoder1989,McDougallMcIntosh2001}.
    
  \end{itemize}
  
  
\end{frame}

 
\section{Mesoscale eddy diffusion and advection in a stratified ocean}

\begin{frame}{Neutral and dianeutral orientation of tracer fluxes}

  Ocean tracer transport is characterized according to whether it is
  dominantly dianeutral (e.g., small scale turbulence) or neutral
  (e.g., mesoscale/geostrophic turbulence).

\begin{figure}  
\centering
\includegraphics[width=.5\textwidth]{FIGS/small_scale_diffusion.pdf}
\caption{\scriptsize Schematic from \textcite{McDougall_etal_geometry}.}
\end{figure}

\begin{equation}
  \footnotesize 
  \text{neutral slope} =
       {\bm S} = 
     -\left[ 
      \frac{-\alpha \, \nablah \Theta + \beta \, \nablah S}
      {\alpha \, \partial_{z} \Theta + \beta \, \partial_{z}S }
      \right]
     \quad 
     \alpha = -\frac{1}{\rho}\frac{\partial \rho}{\partial \Theta}
     \quad
     \beta = \frac{1}{\rho}\frac{\partial \rho}{\partial S}.    
\end{equation}



\begin{itemize}
\item[$\star$] Dianeutral = transport across buoyancy surfaces
  \begin{itemize}
    \footnotesize
  \item[$\bullet$] Commonly approximated by vertical diffusion. 
  \end{itemize}
  
\item[$\star$] Neutral = transport along buoyancy surfaces (neutral
  directions)
   \begin{itemize}
     \footnotesize
   \item[$\bullet$] Rotation of the diffusive fluxes to align with
    neutral directions.
  \end{itemize}
  \item[$\star$] For large-scale transport parameterized by diffusion,
   ratio of dianeutral to neutral diffusivities  is
   $10^{-8}$ in ocean interior.
\end{itemize}

      
\end{frame}



\begin{frame}{Key paradigm: neutral diffusion + secondary circulation}

 \begin{itemize}
      \item[$\star$] Downscale turbulent cascade of tracer variance
        oriented by buoyancy leads to the notion of {\bf neutral
          diffusion}.
      \begin{itemize}
      \scriptsize
      \item[$\bullet$] \textcite{Solomon1971}, \textcite{Redi1982},
        \textcite{GriffiesIsoneutral},
        \textcite{McDougall_etal_geometry}
      \end{itemize}
        
    \item[$\star$] Baroclinic eddies reduce available potential energy
      (restratify) via a secondary overturning circulation. Effects on
      tracers are parameterized by an {\bf eddy advection} or {\bf
        skew diffusion}.
     \begin{itemize}
      \scriptsize
    \item[$\bullet$] \textcite{GM90}, \textcite{GM95},
      \textcite{GriffiesSkew}
      \end{itemize}
    
    \end{itemize}

\begin{figure}
\centering
\includegraphics[width=.5\textwidth]{FIGS/GentMcWilliamsEffect_and_neutral_diffusion.pdf}
\caption{\scriptsize Schematic of the neutral diffusion process
  (mixing along neutral directions) and the Gent-McWilliams effect
  (restratification of density via a secondary circulation that
  reduces available potential energy).}
\end{figure}

\end{frame}




\section{Neutral diffusion: parameterization of eddy-induced tracer mixing}


\begin{frame}{\textcite{Veronis1975}: serious problems with geopotential diffusion}

  \begin{itemize}
    \footnotesize
    \item[$\star$] Geopotential oriented diffusion is problematic
      particularly in regions of steep neutral directions (e.g., western boundary currents like Gulf Stream).
     \item[$\star$] Spurious dianeutral diffusivity scales like
       $A^{\mathsf{\scriptscriptstyle hordif}} \, |{\bm S}|^{2}
       \approx (10^{-3}$ to $10^{-1})~\mbox{m}^{2}~\mbox{s}^{-1}$ with
       slope $|{\bm S}| =10^{-3}$ to $10^{-2}$ and
       $A^{\mathsf{\scriptscriptstyle hordif}} \approx
       10^{3}~\mbox{m}^{2}~\mbox{s}^{-1}$.
     \item[$\star$] Western bdy upwelling in the
       \textcite{Holland1971} simulation is due to spurious dianeutral
       mixing from horizontal diffusion:
       $w \, \partial_{z} \rho = A^{\mathsf{\scriptscriptstyle hordif}}
       \, \partial_{xx}\rho$.

  \item[$\star$] This unphysical circulation is known as the "Veronis
    Effect".
    %\footnote{\scriptsize Veronis acknowledged Henry Stommel
    %for the key insight: "{\it In a discussion about the Holland
    %  (1971) model, H. Stommel pointed out to me the effective
    %  vertical [dianeutral] diffusion associated w/ horizontal mixing
    %  along level surfaces.}"}
    
  \end{itemize}

\begin{figure}
\centering
\includegraphics[width=.75\textwidth]{FIGS/rm98_figure2.pdf}
\caption{\scriptsize \textcite{RobertsMarshall1998} schematic of the
  Veronis Effect.}
\end{figure}


\end{frame}



\begin{frame}{Neutral diffusion}

\begin{itemize}
    \item[$\star$] \textcite{Redi1982} proposed the following neutral
      diffusion flux
      \begin{equation}
        {\bm F}^{\mathsf{\scriptscriptstyle redi}} =
       -\mathbb{K}^{\mathsf{\scriptscriptstyle redi}} \cdot \nabla C =
      -\frac{A^{\mathsf{\scriptscriptstyle ndif}} }{1+S^{2}} 
 \left[ \begin{array}{ccc}
 1 + S_{y}^{2} 
 &
 -S_{x}S_{y}
 &
 S_{x}
 \\
 -S_{x}S_{y}
 &
 1 + S_{x}^{2} 
 &
 S_{y}
 \\
 S_{x}
 &
 S_{y}
 &
 S^{2}
\end{array} 
   \right]
  \left[
        \begin{array}{c}
          \partial_{x}C  \\
          \partial_{y}C  \\
          \partial_{z}C
      \end{array}
        \right]
\label{eq:redi-full-tensor}
\end{equation}

\item[$\star$] Projection operator form \parencite{Olbersetal1985} is
  more intuitive:
\begin{equation}
    \mathbb{K}^{\mathsf{\scriptscriptstyle redi}}
    = A^{\mathsf{\scriptscriptstyle ndif}} \, (\mathbb{I}  - \hat{\gamma} \otimes \hat{\gamma})
    \qquad 
    \hat{\bm \gamma} = 
     \frac{-\alpha \nabla \Theta + \beta \nabla S}
          {|-\alpha \nabla \Theta + \beta \nabla S|}.
\end{equation}

\begin{figure}
\centering
\includegraphics[width=.5\textwidth]{FIGS/neutral_diffusion.pdf}
\caption{\scriptsize Schematic of the neutral diffusion process. The
  downgradient neutral diffusive fluxes are oriented along neutral
  directions.}
\end{figure}

\end{itemize}



\end{frame}


\begin{frame}{Small slope neutral diffusion and boundary regularization}

  \begin{itemize}
    \scriptsize 
    \item[$\star$] \textcite{Solomon1971} and \textcite{GM90} wrote
      down the (small neutral slope) diffusive tracer flux, which is
      used by modelers:
      \begin{equation}
        \small 
      {\bm F} =
       -\mathbb{K}^{\mathsf{\scriptscriptstyle small}} \cdot \nabla C =
      -A^{\mathsf{\scriptscriptstyle ndif}}
      \left[
        \begin{array}{ccc}
          1 & 0 & S_{x} \\
          0 & 1 & S_{y} \\
          S_{x} & S_{y} & 0
      \end{array}
        \right]
      \left[
        \begin{array}{c}
          \partial_{x}C  \\
          \partial_{y}C  \\
          \partial_{z}C
      \end{array}
        \right].
\end{equation}

 \item[$\bullet$] Regularization is required near boundaries and
   regions of steep neutral slopes:
   \textcite{Cox1987,GKW91,DanaMcWilliams1995,ferrari_mcwilliams2008,dana_ferrari_mcwilliams2008,ferrari_griffies_nurser_vallis2010}.
     
\end{itemize}

\begin{figure}
\centering
\includegraphics[width=.6\textwidth]{FIGS/Ferrari_etal2008.pdf}
\caption{\scriptsize Neutral diffusion near surface boundary from
  \textcite{ferrari_mcwilliams2008}. As eddies encounter the boundary,
  they can only mix horizontally rather than along neutral directions.
  Regularization involves tapering to zero the off-diagonal terms in
  the diffusion tensor. Regularization manifests physics while
  supporting numerical stability.}
\end{figure}


\end{frame}


\begin{frame}{Numerical/physical details are crucial}

\begin{itemize}
         \footnotesize
         \item[$\star$] Potential for a nonlinear instability when
           diffusing temperature and salinity: nonlinear diffusion can
           rip the density surfaces apart if not careful.

         \item[$\star$] Nonlinear numerical instability is removed if
           the numerical realization satisfies the neutrality
           condition: $\alpha \, {\bm F}(\Theta) = \beta \, {\bm
             F}(S)$ \parencite{GriffiesIsoneutral}.
           
%       \item [$\star$] We generally cannot ensure positive-definite
%         tracer concentration with local fluxes:
%         \textcite{GriffiesIsoneutral}, \textcite{Beckersetal1998},
%         \textcite{Gnanadesikan1999c}, \textcite{Lemarie_etal2012b},
%         \textcite{Groeskamp_slopes2019}.
         
%         \item [$\star$] Yet non-local fluxes in the
%           \textcite{Shao_etal_2020} scheme ensure
%           positive-definiteness of the rotated diffusion.
      
\end{itemize}

\begin{figure}
\centering
\includegraphics[width=.8\textwidth]{FIGS/GGPLDS_fig6_9.pdf}
\caption{\scriptsize From \textcite{GriffiesIsoneutral}: left panel
  shows triad stencil for neutral diffusive fluxes arising from a
  functional discretization; right panel illustrates the nonlinear
  numerical instability realized by the \textcite{Cox1987} method and
  remedied by \textcite{GriffiesIsoneutral}.}
\end{figure}


\end{frame}


\begin{frame}{Refinements to numerical neutral diffusion over the years}

\begin{itemize}

\item[$\star$] \textcite{Beckersetal1998,Beckers2000} advanced a
   no-go theorem, stating that any linear, consistent, local
   implementation of rotated diffusion can produce extrema.
   \begin{itemize}
     \scriptsize 
     \item[$\bullet$] There are no locally downgradient rotated
       diffusive fluxes on a lattice.
   \end{itemize}
   
 \item[$\star$] \textcite{Lemarie_etal2012a,Lemarie_etal2012b}
   suggested a biasing of \textcite{GriffiesIsoneutral} triad scheme
   to reduce extrema.

 \item[$\star$] \textcite{Groeskamp_slopes2019,Shao_etal_2020}
   developed a non-local method that ensures local downgradient
   properties and thus removes extrema.

\end{itemize}

\begin{figure}
  \centering
  \includegraphics[width=.7\textwidth]{FIGS/lemarie.pdf}
\end{figure}

\end{frame}



\section{Gent-McWilliams effect: parameterization of eddy induced restratification}


\begin{frame}{The \textcite{GM95} eddy parameterization}

\begin{itemize}

\item[$\star$] \textcite{GM90,GM95} noted that mesoscale eddies do
  more than diffuse!
   \begin{itemize}
     \scriptsize 
     \item[$\bullet$] Eddies support a secondary overturning
       circulation that restratifies the flow as it reduces the
       available potential energy.
     \item[$\bullet$] The overturning circulation retains the census
       of water properties within buoyancy layers, rather than mixing
       across layers (``adiabatic'' in ocean physics jargon).
   \end{itemize}
   
 \item[$\star$] \textcite{GM95} proposed a closure for the circulation
   in terms of an eddy-induced streamfunction
   \begin{equation}
     {\bm \Psi} = A^{\mathsf{\scriptscriptstyle gm}} \, \hat{\bm z} \times {\bm S}
     \qquad
         {\bm v}^{*} = \nabla \times {\bm \Psi}
         = -\partial_{z}(A^{\mathsf{\scriptscriptstyle gm}} \, {\bm S})
         \  + \hat{\bm z} \, \nabla \cdot (A^{\mathsf{\scriptscriptstyle gm}} \, {\bm S}).
   \end{equation}
   
\end{itemize}

\begin{figure}
\includegraphics[width=.55\textwidth]{FIGS/GentMcWilliamsEffect.pdf}
\caption{Secondary overturning circulation, ${\bm \Psi}$, arising from
  the \textcite{GM95} eddy-induced parameterization for a southern
  hemisphere front.  This circulation acts to restratify the density,
  thus reducing the available potential energy.}
\end{figure}

\end{frame}


\begin{frame}{Eddy advection and skew diffusion}

\begin{itemize}
\item[$\star$] Since the eddy advection velocity is non-divergent,
  $\nabla \cdot {\bm v}^{*} = 0$, we can realize the tracer transport
  either via an advection or a skew diffusion
  \parencite{GriffiesSkew}.
  \begin{subequations}
    \scriptsize 
        \begin{align}
          {\bm F}^{\mathsf{\scriptscriptstyle adv}}
          &= (\nabla \times {\bm \Psi}) \, C, \;
          {\bm F}^{\mathsf{\scriptscriptstyle skew}}
            = -\nabla C \times {\bm \Psi} 
          \Longrightarrow 
            \nabla \cdot {\bm F}^{\mathsf{\scriptscriptstyle adv}}
            =
            \nabla \cdot {\bm F}^{\mathsf{\scriptscriptstyle skew}}
            \\
 {\bm F}^{\mathsf{\scriptscriptstyle skew}} &=
       -\mathbb{A}^{\mathsf{\scriptscriptstyle skew}} \cdot \nabla C =
      -A^{\mathsf{\scriptscriptstyle gm}}
      \left[
        \begin{array}{ccc}
          0 & 0 & -S_{x} \\
          0 & 0 & -S_{y} \\
          S_{x} & S_{y} & 0
      \end{array}
        \right]
      \left[
        \begin{array}{c}
          \partial_{x}C  \\
          \partial_{y}C  \\
          \partial_{z}C
      \end{array}
        \right].
 \end{align}
\end{subequations}


\end{itemize}

\begin{figure}
\includegraphics[width=.75\textwidth]{FIGS/GentMcWilliamsEffect.pdf}
\caption{Orientation of the \textcite{GM95} circulation, ${\bm \Psi}$,
  and \textcite{GriffiesSkew} density skew flux, ${\bm
    F}^{\mathsf{\scriptscriptstyle skew}}$ (downgradient horizontal;
  upgradient vertical).  Both the circulation and the skew flux
  flatten the density surfaces, thus reducing available potential
  energy.}
\end{figure}

\end{frame}


\begin{frame}{Combination of neutral diffusion plus skew diffusion}

\begin{itemize}

\item[$\star$] Combining the \cite{GM95} scheme as a
  \textcite{GriffiesSkew} skew-diffusion, plus including small angle
  neutral diffusion, yields the subgrid scale transport tensor
  \begin{multline}
 \mathbb{A}^{\mathsf{\scriptscriptstyle skew}} + \mathbb{K}^{\mathsf{\scriptscriptstyle small}}
 =
 \\
 \left[
  \begin{array}{ccc}
 A^{\mathsf{\scriptscriptstyle ndif}} & 0 & (A^{\mathsf{\scriptscriptstyle ndif}} - A^{\mathsf{\scriptscriptstyle gm}}) \, S_{x}
 \\
 0 &  A^{\mathsf{\scriptscriptstyle ndif}} & (A^{\mathsf{\scriptscriptstyle ndif}} - A^{\mathsf{\scriptscriptstyle gm}}) \, S_{y}
 \\
 (A^{\mathsf{\scriptscriptstyle ndif}} + A^{\mathsf{\scriptscriptstyle gm}}) \, S_{x}
 &
 (A^{\mathsf{\scriptscriptstyle ndif}} + A^{\mathsf{\scriptscriptstyle gm}}) \, S_{y}
 &
 A^{\mathsf{\scriptscriptstyle ndif}} \, |{\bm S}|^{2}
 \end{array} 
 \right].
 \label{eq:diffusion-plus-gm-skew}
\end{multline}

\item[$\star$] \textcite{DukSmith1997} suggested that
  $A^{\mathsf{\scriptscriptstyle ndif}} =
  A^{\mathsf{\scriptscriptstyle gm}}$.
  \begin{itemize}
      \footnotesize 
     \item[$\bullet$] Horizontal tracer fluxes reduce to downgradient
       horizontal diffusion.
       \item[$\bullet$] Greatly simplifies the model implementation.
   \end{itemize} 

\item[$\star$] However, further research pointed to fundamental
  distinctions.
  \begin{itemize}
    \footnotesize
   \item[$\bullet$] Eddy stirring that leads to neutral diffusive
     mixing is a distinct process from eddy-induced restratification.
   \item[$\bullet$] Distinct boundary conditions:
     \textcite{McDougallMcIntosh2001,ferrari_mcwilliams2008,ferrari_griffies_nurser_vallis2010}.
   \item[$\bullet$] Distinct diffusivities:
     \textcite{SmithMarshall2009,Abernathey2013}.
   \end{itemize} 

\end{itemize}
\end{frame}


\begin{frame}{Isopycnal averaging, bolus velocity, and Stokes' drift}

\begin{figure}
\includegraphics[width=.9\textwidth]{FIGS/stokes_drift.pdf}
\caption{A single layer of constant density fluid, with resting
  thickness $h=\ho$ and instantaneous thickness $h=\ho + h'(y,t)$.
  Sinusoidal undulations in thickness correlate with fluctuations in
  the meridional velocity, $v' = \vo \, \sin(k \, y - \omega \, t)$}
\end{figure}

\begin{itemize}
  
\item[$\star$] Velocity fluctuates to the right ($v' > 0$) under a
    positive thickness anomaly ($h' > 0$) and to the left ($v' < 0$)
    with a negative anomaly ($h' < 0$).

\item[$\star$] The nonzero correlation between $h'$ and $v'$ leads to
  a Stokes drift, with $\overline{v' \, h'}/\overline{h}$ referred to
  as the bolus velocity \parencite{Rhines1982}.

\item [$\star$] This of the peristaltic contractions and expansions
    that move food through the digestive system.

\end{itemize}

\end{frame}



\begin{frame}{Eddy form stress and thermal wind shears}

 In a geostrophic flow, horizontal density gradients give rise to
 vertical {\it thermal wind} shears in the horizontal flow
\begin{equation}
  f \, \partial_{z}{\bm u} = -(g/\rhoo) \hat{\bm z} \times \nabla \rho
  = -N^{2} \, \hat{\bm z} \times {\bm S}  
  \quad \mbox{with $f = 2 \, \Omega \, \sin(\text{lat})$.}  
\end{equation}

\vspace{-.25cm} 
\begin{figure}
\includegraphics[width=.45\textwidth]{FIGS/thermal_wind_schematic.pdf}
\caption{Example of thermal wind shear as per the atmospheric jet
  stream in the northern hemisphere.  The density increases to the
  north and downward, supporting an eastward vertical shear of the
  zonal flow.}
\end{figure}

\begin{itemize}
  \scriptsize
  
\item[$\star$] \textcite{GreatbatchLamb1990}: vertical viscous
  friction weakens the thermal wind shear, and in so doing reduces
  baroclinicity: just like \textcite{GM95}!

\item[$\star$] Vertical friction is a simple parameterization of
  vertical eddy form stress:
  \textcite{FerreiraMarshall2006,ZhaoVallis2008,Young2012,Loose_etal2023}.

\end{itemize}

\end{frame}



\section{Closing remarks and further research}


\begin{frame}{Summarizing mesoscale eddy restratification + diffusion}

\begin{itemize}

\item[$\star$] \textcite{GM95} eddy-advection + eddy neutral diffusion: mesoscale eddy
  parameterization appears just in tracer equation
  \begin{equation}
    \partial_{t} C + \nabla \cdot [({\bm v} + {\bm v}^{*}) \, C] = \nabla \cdot (\mathbb{K} \cdot \nabla C).
  \end{equation}
  
\item[$\star$] \textcite{GriffiesSkew} eddy skew-diffusion + eddy neutral 
  diffusion: mesoscale eddy parameterization appears just in tracer
  equation
  \begin{equation}
    \partial_{t} C + \nabla \cdot ({\bm v} \, C) =
      \nabla \cdot [ (\mathbb{K} + \mathbb{A}^{\mathsf{\scriptscriptstyle skew}})  \cdot \nabla C].
  \end{equation}
  
\item[$\star$] \textcite{GreatbatchLamb1990} + eddy neutral diffusion:
  parameterized eddies in tracer + velocity equations
  \begin{subequations}
    \begin{align}
    \partial_{t} C + \nabla \cdot ({\bm v} \, C) &= \nabla \cdot (\mathbb{K} \cdot \nabla C)
      \\
      \partial_{t} {\bm u} + ({\bm v} \cdot \nabla) {\bm u}
      + f \, \hat{\bm z} \times {\bm u} &=
      -(1/\rhoo) \, \nablah p + \partial_{z}(\kappa^{\mathsf{\scriptscriptstyle eddy}} \, \partial_{z}{\bm u}).
    \end{align}
  \end{subequations}
  
\end{itemize}

\end{frame}


\begin{frame}{Further related research}


Research into mesoscale eddy parameterizations remains an active and
challenging area.  In addition to those topics mentioned earlier, here
are a few others. Warning: this is an incomplete list!

\begin{itemize}

  \scriptsize
  
\item[$\star$] {\bf Anisotropic diffusion/advection}:
  \textcite{SmithGent2004,Foxkemper_etal2013,bachman2020diagnosis}
  pursued laterally anisotropic approaches, with mixing oriented
  according to the flow.

\item[$\star$] {\bf Relating neutral diffusion and eddy advection}:
  \textcite{SmithMarshall2009,Abernathey2013}

\item[$\star$] {\bf Energy informed closure}:
  \textcite{EdenGreat2008,MarshallAdcroft2010,Grooms_etal2025}

\item[$\star$] {\bf Non-Newtonian stresses}:
  \textcite{AnsteyZanna2017,Bachman_etal2018}

\item[$\star$] {\bf Vorticity/energy informed closure}:
  \textcite{Marshall_etal2012,Maddison_Marshall2013,Mak_etal2017}

\item[$\star$] {\bf Backscatter for the eddy permitting regime}:
  \textcite{Jansen_Held_2014,Jansen_etal2015,Juricke_etal2020,Pudig_etal2026}

\item[$\star$] {\bf Vertical structure of the diffusivity}:
  \textcite{zhang2024role,Yankovsky_etal2024}
  

\end{itemize}  

\end{frame}




\begin{frame}{Some pointers for early career researchers}

\begin{itemize}

 \item[$\star$] Developing subgrid scale parameterizations requires a
   synergy between physics, maths, and numerics.
  \begin{itemize}
  \scriptsize 
  \item[$\bullet$] No one approach is sufficient nor more or less important.
  \item[$\bullet$] Collaboration is needed, though with a strong dose
    of individual creative (and obsessive) passion to open up new
    avenues that break from the crowd.
  \end{itemize}

\item[$\star$] Read papers and books! 
  \begin{itemize}
    \scriptsize
    \item[$\bullet$] Do not rely on AI to summarize the literature. 
  \item[$\bullet$] It can be enlightening to see how previous
    researchers took wrong turns, only to eventually find their way.
  \item[$\bullet$] Reading is the best training for writing.
  \end{itemize}

 \item[$\star$] Patience and persistence generally lead to progress.    
   \begin{itemize}
     \scriptsize
     \item[$\bullet$] No theory nor method is beyond being questioned
       or overturned.     
     \item[$\bullet$] But be ready for resistence. Many papers that
       broke new ground were initially rejected or resisted (sometimes
       for many years).
  \item[$\bullet$] ``Slow science'' can lead to the biggest
    breakthroughs and longer lasting results.
  \end{itemize}    
  
\end{itemize}

\end{frame}


\begin{frame}{Un grand merci aux organisateurs pour leur accueil \`{a} Paris !}

\begin{figure}
\centering
\includegraphics[width=.9\textwidth]{FIGS/CM2.6_SPEED_ICE_2000528_NA_MILLER.png}
\caption{\scriptsize Ocean surface speed from a global climate
 model (GFDL/CM2.6) with roughly 10~km horizontal grid spacing in
 the middle-Atlantic Ocean \parencite{Griffies_etal_CM2p6_heat}.}
\label{fig:CM2p6_NATl_speed}
\end{figure}



\end{frame}

\section{References} 

\begin{frame}[allowframebreaks]{References}
    \printbibliography
\end{frame}



\section{Extra slides}


\begin{frame}{Buoyancy and associated watermasses orient ocean flows}


\begin{figure}
\centering
\includegraphics[width=.475\textwidth]{FIGS/WunschFerrariARFM2004_fig1.pdf}
\hspace{.1cm}
\includegraphics[width=.475\textwidth]{FIGS/Sverdrup_schematic.pdf}
\caption{\scriptsize Left: neutral density section in Pacific, from
  \textcite{Wunsch_Ferrari2004}. Right: Schematic of Southern Ocean
  circulation from \textcite{SverdrupJohnsonFleming1942}. Note: south
  is swapped in these two depictions.}
\end{figure}

Physical processes (e.g., waves, eddies, mixing, circulation) are
oriented according to buoyancy.
      \begin{itemize}
      \scriptsize
      \item[$\bullet$] \textcite{Montgomery1938},
        \textcite{Iselin1939}, \textcite{SverdrupJohnsonFleming1942},
        \textcite{Walin1977}, \textcite{McDougall1987a},
        \textcite{McDougall_etal_geometry}, \textcite{Tailleux2026}
      \end{itemize}


\end{frame}



\begin{frame}{The mesoscale part of the ocean physics zoo}

\begin{figure}
\centering
\includegraphics[width=.55\textwidth]{FIGS/Fig1_OceanProcSchem_revision3.png}
\caption{\scriptsize From \textcite{Griffies_Treguier2013}.}
\end{figure}

We here focus on the ocean mesoscale, which is characterized by
     \begin{itemize}
       \footnotesize
     \item[$\bullet$] Small Rossby number (rotationally dominant)
     \item[$\bullet$] Large Richardson number (strong vertical stratification)  
     \item[$\bullet$] Quasi-geostrophy (flow scales $L \approx
       \mathcal{O}(L_{d})$) and planetary geostrophy (flow scales $L
       \gg \mathcal{O}(L_{d})$) inform the parameterizations ($L_{d}$
       is the deformation scale).
    \end{itemize}
    

\end{frame}




\begin{frame}{Neutral directions: \textcite{McDougall1987a,McDougall1987b}}

%\begin{figure}
\begin{center}
\includegraphics[width=.5\textwidth]{FIGS/Encylcopedia_Neutral_figure4.pdf}
%\caption{
\newline 
{\scriptsize Neutral trajectory and $\gamma^{n}$ surface bounded by
  $\sigma_{0}$ and $\sigma_{4}$. Atlantic section.  From
  \textcite{McDougallJackett2009}.}
%}
\end{center}
%\end{figure}

\begin{itemize}
    \item[$\star$] Importance of local referencing to define neutral
      directions for rotating the diffusion tensor.
    
    \item[$\star$] Neutral orientation empirically justified by small
      levels of measured dianeutral mixing in the ocean interior
      \parencite{McDougall_etal_geometry}.
     \item[$\star$] Further refinement from \textcite{Tailleux2026}.  
\end{itemize}

\end{frame}


\begin{frame}{An aside: neutral helicity: \textcite{McDougallJackett1988}}

%\begin{figure}
\begin{center}
\includegraphics[width=.5\textwidth]{FIGS/neutral_helicity.pdf}
%\caption{
\newline 
    {\scriptsize Following a neutral direction around a loop does not
      return to the same point in space, as first noted by
      \textcite{McDougallJackett1988}.}
%}
\end{center}
%\end{figure}

\begin{itemize}
\item[$\star$] Neutral helicity is nonzero for a realistic ocean
  \begin{equation}
   \mathcal{H}^{\gamma}  
   = \beta \, {\cal T} \, \nabla p \cdot (\nabla S \, \times \, \nabla \Theta)
   \quad \mbox{with} \quad {\cal T} = \beta \, \partial_{p} (\alpha/\beta),
  \end{equation}
 where $\mathcal{H}^{\gamma} \ne 0$ produces a helical structure.    
    
\item[$\star$] Further work by \textcite{Bennett2019},
  \textcite{Stanley2019}, \textcite{Stanley_etal2021}, and
  \textcite{Chatelain_etal2026}.

  
\end{itemize}

\end{frame}




\begin{frame}{Griffies et al (1998): remedied problems with Cox scheme}

\begin{itemize}
 \item[$\star$] Respect neutrality (McDougall 1987) eliminates Cox instability 
 \begin{equation}
     \alpha {\bm F}(\Theta) = \beta {\bm F}(S_{A})
 \end{equation}

 
 \item[$\star$] Functional method led to "triad discretization" ensuring the reduction of global variance
 \begin{align}
     2 {\cal F} = \int ({\bm F} \cdot \nabla C) \, \mathrm{d}V
% &= -\frac{1}{2} \int  (\nabla C \cdot \mathbb{K} \cdot \nabla C) \, \mathrm{d}V
     \le 0 
%     \\
    \qquad 
    \delta {\cal F}/\delta C 
     = -(\nabla \cdot {\bm F}) \, \mathrm{d}V.
 \end{align}
 
\item[$\star$] Scheme is stable: requires no horizontal background diffusion. 
 
\begin{center}
\includegraphics[width=.75\textwidth]{FIGS/GGPLDS_fig6_9.pdf}
\end{center}
 
\end{itemize}

\end{frame}


\begin{frame}{Neutral slope tapering: stability does not imply accuracy}

\begin{itemize}

\item[$\star$] Tapering allows us to run the model stably.

 \item[$\star$] But accuracy is degraded when neutral slopes are steeper than the grid aspect ratio (Beckers et al 2000)
 (akin to the representation of topography).

\end{itemize}


\begin{figure}[ht]
\begin{center}
\includegraphics[width=.8\textwidth]{FIGS/neutral_direction.pdf}
\caption{Zonal-vertical grid cell arrangement \& sample neutral directions.}
\label{fig:neutral_direction}
\end{center}
\end{figure}

\end{frame}


\begin{frame}{Layer (2d convergence) vs. level (3d convergence)}

\begin{itemize}

\item[$\star$] Layer models compute isopycnal (or ideally neutral tangent plane) diffusion as a 2d layer-convergence 
($h^{\gamma} = $ thickness)
\begin{equation}
     {\bm F}^{\mbox{\tiny layer}}= - A \, \nabla_{\gamma} C
     \qquad 
      {\cal R}^{\mbox{\tiny layer}} = 
     -\frac{1}{h^{\gamma} }  
     \left[ \nabla_{\gamma} \cdot (h^{\gamma} \,
     {\bm F}^{\mbox{\tiny layer}} ) \right].
\end{equation}

\item[$\star$] Level models compute small slope neutral diffusion as a 3d level-convergence
\begin{equation}
   {\bm F}^{\mbox{\tiny level}} 
  =  -A \, \left[ \nabla_{\gamma}  + \hat{\bm z} \, ({\bm S} \cdot \nabla_{\gamma} ) \right] C  
   \qquad 
   {\cal R}^{\mbox{\tiny level}} 
    = -\nabla_{z} \cdot {\bm F}_{h}^{\mbox{\tiny level}}   
    - \partial_{z} F_{z}^{\mbox{\tiny level}}
\end{equation}

\item[$\star$] Maths shows that the operators are identical
\parencite{SMGbook,McDougall_etal_geometry}  
\begin{equation}
 {\cal R}^{\mbox{\tiny layer}}  = {\cal R}^{\mbox{\tiny level}}.
\end{equation}

\item[$\star$] Hence, in continuum, neutral diffusion formulated
  layer-wise is the same as small slope neutral diffusion formulated
  level-wise.

\item[$\star$] But they are distinct numerically. Layer approach is
  ensured to be downgradient, whereas level approach (rotated
  diffusion) can suffer from extrema ala \textcite{Beckersetal1998}.


\end{itemize}

\end{frame}



\begin{frame}{New method: vertically non-local slope/gradient estimate}

\begin{itemize}

\item[$\star$] \textcite{Groeskamp_slopes2019}: observational analysis
  method focused on cabbeling and ventilation diagnostics (level 3d
  convergence implementation).

\item[$\star$] \textcite{Shao_etal_2020}: MOM6 ALE-layer neutral
  diffusion (2d convergence to ensure no extrema).

\item[$\star$] Slope tapering imposed by grid resolution: cannot
  represent a neutral direction so steep that it leaves domain within
  a cell.

\item[$\star$] Hence, "max slope" is chosen by grid resolution not
  modeler.


\end{itemize}

\begin{center}
\includegraphics[width=.5\textwidth]{FIGS/groeskamp.pdf}
\label{fig:groeskamp}
\end{center}

\end{frame}



\begin{frame}{Neutral skew diffusion from polarized tracer fluxes}

\begin{itemize}

\item[$\star$] Layer averaged tracer equation (\textcite{Smith1999},
  Section 9.4.3 of \textcite{SMGbook})
\begin{align}
 (\partial_{t} + \widehat{\bm u} \cdot \nabla_{\gamma} ) \,  \widehat{C}
 &= -\frac{1}{\overline{h}^{\gamma}} \, 
   \nabla_{\gamma} \cdot (\overline{h}^{\gamma} 
   \, \widehat{C'' \, {\bm u}''})
  \\
  \widehat{\Phi} = \frac{\overline{\Phi \, h}^{\gamma}}{\overline{h}^{\gamma}}
  \qquad \Phi &= \widehat{\Phi} + \Phi'' 
\qquad
\widehat{\bm u} = \overline{\bm u}^{\gamma} + {\bm u}^{\mbox{\tiny bolus}}
  \\
  \widehat{C'' \, {\bm u}''} 
  &= - \mathbb{M} \cdot \nabla_{\gamma} \widehat{C}
\label{eq:thickness-tracer-evolution}
\end{align}

\item[$\star$] The $\mathbb{M}$ tensor arises from correlations
  between tracer and velocity on the neutral tangent plane.  It has
  both a symmetric and an anti-symmetric component.

\item[$\star$] Polarized tracer fluxes \parencite{MiddletonLoder1989}
  contribute to the anti-symmetric component of $\mathbb{M}$.

\item[$\star$] We presently only parameterize the symmetric component
  of $\mathbb{M}$ via neutral diffusion.

\item[$\star$] What justifies ignoring unresolved polarized tracer
  fluxes?
  
\end{itemize}


\end{frame}




\end{document}
